\labeledsection{Functional Programming using SML}{sec:sml}
\definition{Functional Programming}{
We call programming languages where procedures can be fully describes in terms of their input/output behabior \textbf{functional}.
}{functional_programming}

\commandnote{
There are two main programming paradigms: \textbf{imperative} and \textbf{declarative}. Imperative programming is where the programmer instructs the machine how to change its state, it can be:
\begin{itemize}
    \item structured programming: makes extensive use of block structures and subroutines
    \item procedural programming: groups instructions into procedures
    \item object-oriented programming: groups instructions together with the part of the state they operate on
\end{itemize}

On the other hands, declarative programming is where the programmer merely declares properties of the desired result, but not how to compute it. Examples:
\begin{itemize}
    \item functional programming: the desired result is declared as the value of a series of function calls
    \item logic programming: the desired result is declared as the answer to a question about a system of facts and rules
    \item mathematical programming: the desired result is declared as the solution of an optimization problem
\end{itemize}
}

\definition{Standard ML (SML)}{
SML is general-purpose, high-level, modular, \linkterm{function programming language}{functional_programming} with compile-time type checking and type inference.
}{def:sml}

\definition{Predefined objects in SML}{
\begin{itemize}
\item \textbf{basic types}, e.g. \texttt{int}, \texttt{real}, \texttt{bool}, \texttt{string}, ...
\item \textbf{basic type constructors}, e.g. \texttt{->}, \texttt{*}, ...
\item \textbf{basic operators}, e.g. \texttt{numbers}, \texttt{true}, \texttt{false}, \texttt{+}, \texttt{-}, \texttt{*}, \texttt{-}, \texttt{>}, \^{} ...
\item \textbf{control structures}, e.g. \texttt{if} $\varPhi$ \texttt{then} $E_1$ \texttt{else} $E_2$;
\item \textbf{comments}, e.g. \texttt{(*  this is a comment *)}
\end{itemize}
}{predefined_objects_sml}

A simple example, if we write \texttt{2;} we get back \texttt{val it = 2 : int}. and if we write \texttt{4.711;} we get back \texttt{val it = 4.711 : real}. In other words, \linkterm{SML}{def:sml} interpreter has determined that the input has a value and that value has a specific \linkterm{type}{predefined_objects_sml}.

\definition{Declarations bind variables}{
\begin{itemize}
    \item \textbf{value declarations}, e.g. \texttt{val pi = 3.1415;}
    \item \textbf{type declarations}, e.g. \texttt{type twovec = int * int;}
    \item \textbf{function declarations}, e.g. \texttt{fun square (x:real) = x * x;}
\end{itemize}
A \linkterm{function declaration}{sml_declarations} only \linkterm{declares}{sml_declarations} the function name as a globally visible name. The \textbf{formal parameters} in brackets are only visible in the \textbf{function body}.
}{sml_declarations}

Alternatively, we can just \linkterm{declare}{sml_declarations} a function object, e.g. via \texttt{val foo = fn x => x + x;}.

\linkterm{SML}{def:sml} \linkterm{functions}{function} that have been \linkterm{declared}{sml_declarations} can be applied to arguments of the right type, e.g. \texttt{square 4.0} which evaluates to \texttt{4.0 * 4.0} and this to \texttt{16.0}.

\inlinecode{- fun square(x:real) = x * x; (*function declaration*)}

\inlinecode{val square = fn : real -> real (*interpreter response*)}

\inlinecode{- square 4.0; (*function application*)}

\inlinecode{val it = 16 : real (*interpreter response*)}

\definition{Local Declaration}{
A local declaration uses \texttt{let} to bind variables in its scope (delineated by \texttt{in} and \texttt{end}).
}{local_declaration_sml}

\commandnote{
\linkterm{Local declarations}{local_declaration_sml} can \textbf{shadow} existing variables. i.e. it hides the old binding in its scope.
}

Example:

\inlinecode{- val test = 4;}

\inlinecode{val test = 4 : int}

\inlinecode{- let val test = 7 in test * test end;}

\inlinecode{val it = 49 : int}

\inlinecode{- test;}

\inlinecode{val it = 4 : int}

\commandnote{
In imperative programming, variable are identifiers for memory locations that can be changed at will. However, in \linkterm{functional programming}{functional_programming}, variables are bound to \linkterm{values}{sml_declarations} in a given \linkterm{scope}{scope}, that \linkterm{value}{sml_declarations} cannot be changed in the \linkterm{scope}{scope}. i.e. variable binding in \linkterm{functional programming}{functional_programming} is a dynamic mapping from variable names to immutable memory locations, which can be locally overwritten (shadowing).
}

\definition{Pattern Matching}{
Pattern matching is using special expressions with unbound variables (we call them structured patterns) to \linkterm{declare}{sml_declarations} more than one variable.
}{pattern_matching_sml}

For example:

\inlinecode{- val unitvector = (1,1);}

\inlinecode{val unitvector = (1,1): int * int}

\inlinecode{- val (x,y) = unitvector;}

\inlinecode{val x = 1 : int}

\inlinecode{val y = 1 : int}

\definition{Anonymous Variables}{
Variables without a name, denoted by "underscore" (\texttt{\_}).
}{anonymous_variables_sml}
e.g. if we are not interested in one if the values: \inlinecode{- val (x,\_) = unitvector} $\leadsto$ \inlinecode{val x = 1 : int}. Similarly, we can define a function that takes a pair and select the first element:

\inlinecode{- fun first(p) = let val (x,\_) = p in x end;} $\leadsto$ \inlinecode{val first = fn : 'a * 'b -> 'a}. Example usage:
\begin{verbatim}
- first (3,"hello");
val it = 3 : int
\end{verbatim}

A selector of the second element:
\begin{verbatim}
- fun second(p) = let val (_,x) = p in x end;
val second = fn : 'a * 'b -> 'b
- second (3,"Welcome to SML");
val it = "Welcome to SML" : string
\end{verbatim}

\commandnote{
\texttt{'a}, \texttt{'b}, ... are called type variables (\linkterm{SML}{def:sml} supports universal types). 
}
So the \texttt{first} \linkterm{function}{function} signature (\texttt{fn : ’a * ’b -> ’a}) means the \linkterm{function}{function} takes a pair of type \texttt{ ’a * ’b} as input and returns an object of type  \texttt{’a}.

\linkterm{SML}{def:sml} has a built-in "list type". Given a type $\langle \langle$ \texttt{ty} $\rangle \rangle$, \texttt{list }$\langle \langle$ \texttt{ty} $\rangle \rangle$ is also a type (a list of objects of type $\langle \langle$ \texttt{ty} $\rangle \rangle$). For example, \texttt{list int} is a type for lists of integers. 
\begin{verbatim}
- [1,2,3];
val it = [1,2,3] : int list
\end{verbatim}

\definition{Lists Constructors}{
Lists have two constructors: \texttt{nil} (empty list) and \texttt{::} (called \texttt{cons}).
}{lists_constructors_sml}
\begin{verbatim}
- nil;
val it = [] : 'a list
- 9::nil;
val it = [9] : int list
\end{verbatim}

\definition{Recursion}{
We call an \linkterm{SML}{def:sml} \linkterm{function}{sml_declarations} recursive, iff the \linkterm{function}{sml_declarations} is called in the body of the \linkterm{function declaration}{sml_declarations}.
}{recursion}

\linkterm{Recursion}{recursion} example:
\begin{verbatim}
- fun upto(m,n) = if m>n then nil else m::upto(m+1,n);
vap upto = fn : int * int -> int list
upto (2,5);
val it = [2,3,4,5] : int list
\end{verbatim}

Let's trace that \linkterm{function}{function}: 
\begin{verbatim}
upto (2,5) = if 2>5 then nil else 2::upto(3,5)
upto(3,5) = if 3>5 then nil else 3::upto(4,5)
upto(4,5) = if 4>5 then nil else 4::upto(5,5)
upto(5,5) = if 5>5 then nil else 5::upto(6,5)
upto(6,5) = if 6>5 then nil else 6::upto(7,5)
\end{verbatim}

so it will evaluate to: \inlinecode{2::3::4::5::nil}

\textbf{Another Example}, consider we want a \linkterm{function}{sml_declarations} that sums all the elements in a list. We also do that via \linkterm{recursion}{recursion}, if we get an empty list, we return 0, if we get a list with elements (a head element, and tail of the rest) then we sum the head to the sum of the tail.

\begin{verbatim}
- fun sumlist nil = 0 | sumlist (h::t) = h + sumlist t;
val sumlist = fn : int list -> int
- sumlist ([1,4,6]);
val it = 11 : int
- sumlist ([]);
val it = 0 : int
\end{verbatim}

\commandnote{
Every \textit{infix} operator in \linkterm{SML}{def:sml} (+,-,...) has a \textit{function style} (\texttt{op+}, \texttt{op-},...) notation.
}
For example, we can write the summation using that notation in our \texttt{sumlist} \linkterm{function}{sml_declarations}:
\begin{verbatim}
- fun sumlist nil = 0 | sumlist (h::t) = op+ (h, sumlist t);
val sumlist = fn : int list -> int

- sumlist ([1,2,3]);
val it = 6 : int
\end{verbatim}
\linkterm{SML}{def:sml} have other \textit{function style} operations, e.g.:
\begin{verbatim}
(op-) (7, 4);            (* 3 *)
(op*) (6, 5);            (* 30 *)
(op div) (10, 3);        (* 3 *)
(op mod) (10, 3);        (* 1 *)
(op^) ("hi", "there");   (* "hithere" *)
(op::) (1, [2,3]);       (* [1,2,3] *)
(op@) ([1,2], [3,4]);    (* [1,2,3,4] *)
(op=) (3, 3);            (* true *)
(op=) ("a", "b");        (* false *)
(op<>) (3, 4);           (* true *)
(op<>) (5, 5);           (* false *)
(op<)  (2, 5);           (* true  *)
(op<=) (5, 5);           (* true  *)
(op>)  (7, 3);           (* true  *)
(op>=) (3, 4);           (* false *)
\end{verbatim}



Evaluation in \linkterm{SML}{def:sml} is \textbf{call-by-value} i.e. to whenever we encounter a \linkterm{function}{sml_declarations} applied to arguments, we compute the value of the argument first. We write $t_1 \leadsto t_2 \leadsto \cdots \leadsto t_n$ for tracing the \linkterm{recursive}{recursion} arguments $t_i$ through a \linkterm{recursive}{recursion} computation. For example:
\[
\operatorname{upto}\paren{2,4} \leadsto 2::\operatorname{upto}\paren{3,4} \leadsto 2::(3::\operatorname{upto}\paren{4,4}) \leadsto 2::(3::(4::\operatorname{upto}\paren{5,4})) \leadsto 2::(3::(4::\operatorname{nil})) 
\]
\commandnote{
\linkterm{Recursive functions}{recursion} need not to terminate. For example:
\inlinecode{fun diverges (n) = n + diverges (n+1)} which evaluates to: $\operatorname{diverges}(1) \leadsto 1 + \operatorname{diverges}(2) \leadsto 1 + (2 + \operatorname{diverges}(3)) \leadsto \cdots$
}

\textbf{Examples:}

Length of a list:
\begin{verbatim}
- fun length nil = 0 | length (_::t) = 1 + length t;
val length = fn : 'a list -> int
- length ([1,2,3,4,5]);
val it = 5 : int
\end{verbatim}

Fibonacci:
\begin{verbatim}
- fun fib 0 = 0
= | fib 1 = 1
= | fib n = fib (n-1) + fib (n-2);
val fib = fn : int -> int

- fib 6;
val it = 8 : int
\end{verbatim}


Factorial:
\begin{verbatim}
- fun fact 0 = 1 | fact n = n * fact(n-1);
val fact = fn : int -> int
- fact 3
val it = 6 : int
\end{verbatim}
\textbf{Note} that each \linkterm{recursive}{recursion} call has to remember the multiplication to do later. The final result waits until \inlinecode{fact(0)} is evaluated. However, there is another \linkterm{recursion}{recursion} technique called \textbf{tail recursion} where we accumulate the result as we go from recursive to another using a helper variable (accumulator) and function. For example, here is the same factorial \linkterm{function}{sml_declarations} using tail recursion:
\begin{verbatim}
- fun facthelper (0, acc) = acc | facthelper (k,acc) = facthelper (k-1, k * acc);
val facthelper = fn : int * int -> int
- fun fact n = facthelper (n,1);
val fact = fn : int -> int
- fact 3;
val it = 6 : int
\end{verbatim}
Let's track that:
$\operatorname{fact}(3) \leadsto \operatorname{facthelper}(3,1) \leadsto \operatorname{facthelper}(2, 3) \leadsto \operatorname{facthelper}(1, 6) \leadsto \operatorname{facthelper}(0,6) \leadsto 6$





\definition{Tail Recursion}{
A tail call is a subroutine call performed as the final action of a routine $R$. If the call is a \linkterm{recursive call}{recursion}, $R$ is said to be tail recursive.
}{tail_recursion}

\definition{Append \texttt{@}}{
The operator \texttt{@} appends two lists together, e.g. \texttt{- [1]@[2];} gives us \texttt{val it = [1,2] : int list}
}{append_sml}
But we can define our own append \linkterm{recursively}{recursion}:
\begin{verbatim}
- fun append (nil,y) = y 
    | append (h::t, y) = h::append (t,y);

val append = fn : 'a list * 'a list -> 'a list

- append ([1,2,3], [4,5,6]);
val it = [1,2,3,4,5,6] : int list
\end{verbatim}

\textbf{Reversing a list:}
\begin{verbatim}
- fun reverse nil = nil
  |   reverse (h::t) = reverse t @ [h];

val reverse = fn : 'a list -> 'a list

- reverse ([1,2,3,4,5,6,7]);
val it = [7,6,5,4,3,2,1] : int list
\end{verbatim}

\textbf{Membership Check:}
\begin{verbatim}
- fun member (_, nil) = false
    | member (x, h::t) =
      if x = h then true
      else member (x,t);

val member = fn : ''a * ''a list -> bool

- member (4,[1,2,3,4,5]);
val it = true : bool

- member (4,[1,2,3,5]);
val it = false : bool
\end{verbatim}

\textbf{Reverse a list and append it to another list:}
\begin{verbatim}
- fun reverse (nil,y) = y
    | reverse (h::t, y) = reverse (t, h::y);

val reverse = fn : 'a list * 'a list -> 'a list

- reverse ([4,5,6], [3,2,1]);
val it = [6,5,4,3,2,1] : int list
\end{verbatim}
Let's trace that with \texttt{[2,3], [1]}:
$\operatorname{reverse}([2,3],[1]) \leadsto \operatorname{reverse}([3],[2, 1]) \leadsto \operatorname{reverse}([],[3,2,1]) \leadsto [3,2,1]$

\textbf{Flattening a list of lists to one list}
\begin{verbatim}
- fun flat nil = nil
    | flat (h::t) = h @ flat t;

val flat = fn : 'a list list -> 'a list

- flat ([["when", "shall"], ["we"], ["meet", "again"]]);
val it = ["when","shall","we","meet","again"] : string list
\end{verbatim}

\textbf{Remove First Occurrence:}

Here is a function that takes an element and a list and give the list without the first occurrence of that element:
\begin{verbatim}
- fun remove_first (x,[]) = []
    | remove_first (x, h::t) =
        if x = h then t else h::remove_first (x,t);

val remove_first = fn : ''a * ''a list -> ''a list

- remove_first (1,[2,3,1,4,1]);
val it = [2,3,4,1] : int list
\end{verbatim}

\textbf{Remove All Occurrences:}
\begin{verbatim}
- fun remove_all (x,[]) = []
    | remove_all (x, h::t) =
        if x = h then remove_all (x,t) else h::remove_all(x,t);

val remove_all = fn : ''a * ''a list -> ''a list

- remove_all (1,[2,3,1,4,1]);
val it = [2,3,4] : int list 
\end{verbatim}

\textbf{Count All Occurrences}

\begin{verbatim}
- fun count _ [] = 0
    | count x (h::t) = (if x = h then 1 else 0) + count x t;

val count = fn : ''a -> ''a list -> int

- count 1 [1,2,3,1,4,1,1,5];
val it = 4 : int
\end{verbatim}

\textbf{Sort List:}

\begin{verbatim}
- fun sort [] = []
    | sort (h::t) =
        let
            fun place (x, []) = [x]
              | place (x, h::t) =
                if x <= h then x :: h :: t else h :: place (x, t)
        in
            place (h, sort t)
        end;

val sort = fn : int list -> int list

- sort [4,5,2,3,8,7,6,1,0];
val it = [0,1,2,3,4,5,6,7,8] : int list
\end{verbatim}
Let's trace that on [2,1,3]:

sort [2,1,3] $\leadsto$
place (2, sort [1,3]) $\leadsto$
place (2, place (1, sort [3])) $\leadsto$
place (2, place (1, place (3, sort []))) $\leadsto$

place (2, place (1, place (3, []))) $\leadsto$
place (2, place (1, [3])) $\leadsto$
place (2, [1,3]) $\leadsto$
1::place(2, [3]) $\leadsto$
1::[2,3] $\leadsto$
[1,2,3] 

\textbf{Merge two Sorted Lists:}
\begin{verbatim}
- fun merge ([],b) = b
    | merge (a, []) = a
    | merge (h1::t1, h2::t2) =
        if h1 <= h2 then h1::merge (t1, h2::t2)
        else h2::merge (h1::t1, t2);

val merge = fn : int list * int list -> int list

- merge ([1,3,5], [2,4,6]);
val it = [1,2,3,4,5,6] : int list
\end{verbatim}




\definition{SML Strings}{
\linkterm{SML}{def:sml} treats strings as lists of characters. We can go from \texttt{string} to \texttt{chat list} and vice versa via \texttt{explode} and \texttt{implode}.
}{sml_string}
For example:
\begin{verbatim}
- explode "Ibrahim";
val it = [#"I",#"b",#"r",#"a",#"h",#"i",#"m"] : char list
- implode it;
val it = "Ibrahim" : string
\end{verbatim}

Using that and our reverse \linkterm{function}{sml_declarations}, we can define a \linkterm{function}{sml_declarations} that checks if a \linkterm{string}{sml_string} is \textit{palindrome} or not (read the same from either side, e.g. otto, anna, etc.)
\begin{verbatim}
- fun reverse nil = nil
    | reverse (h::t) = reverse t @ [h];

val reverse = fn : 'a list -> 'a list

- fun palindrome s = 
    let 
        val cs = explode s 
    in 
        cs = reverse cs 
    end;

val palindrome = fn : string -> bool

- palindrome "apa";
val it = true : bool
- palindrome "apaa";
val it = false : bool
\end{verbatim}

\definition{Short-Circuit Operators}{
\linkterm{SML}{def:sml} has built-in short-circuit operators. \texttt{andalso} is the logical AND and \texttt{orelse} is the logical OR. These behave like \texttt{\&\&} and \texttt{||} in C/Java.
}{orelse_andalso}
\begin{verbatim}
- true andalso false;
val it = false : bool

- true orelse false;
val it = true : bool
\end{verbatim}

\definition{Higher-Order Functions}{
In \linkterm{SML}{def:sml}, we call \linkterm{functions}{sml_declarations} that take \linkterm{functions}{sml_declarations} as arguments \textbf{higher-order functions} and those that do not \textbf{first-order functions}
}{higher_order_functions}

For example, the built-in \linkterm{higher-order function}{higher_order_functions} \texttt{map} that takes a \linkterm{functions}{sml_declarations} and a list and returns a list with the passed \linkterm{functions}{sml_declarations} applied to each element of the list.
\begin{verbatim}
- fun f x = x + 1;
val f = fn : int -> int
- map f [1,2,3];
val it = [2,3,4] : int list
\end{verbatim}

\textbf{Note: }we can program our own \texttt{map} \linkterm{function}{sml_declarations} \linkterm{recursively}{recursion}:
\begin{verbatim}
- fun mymap (f, nil) = nil
    | mymap (f, h::t) = f(h)::mymap (f,t);
val mymap = fn : ('a -> 'b) * 'a list -> 'b list

- mymap (f,[1,2,3]);
val it = [2,3,4] : int list
\end{verbatim}

\textbf{Selective Copy:}

We would like to make a function that takes a list and give us copy of the list with, for instance, only the even numbers. First, we define an \texttt{is\_odd} function:
\begin{verbatim}
- fun is_odd x = (x mod 2) = 1;
val is_odd = fn : int -> bool
- is_odd 3;
val it = true : bool
- is_odd 2;
val it = false : bool
\end{verbatim}
Then we define a function \texttt{remove\_if} that takes two arguments, one is the filter (a function as an argument) (here it is our \texttt{is\_odd}), and a list which we hope to filter out. 
\begin{verbatim}
- fun remove_if filter [] = []
    | remove_if filter (h::t) =
        if filter(h) then remove_if filter t
        else h::(remove_if filter t);

val remove_if = fn : ('a -> bool) -> 'a list -> 'a list

- remove_if is_odd [1,2,3,4,5,6,7,8,9];
val it = [2,4,6,8] : int list
\end{verbatim}


\definition{Function Literal}{
A function literal (also called \textbf{anonymous function}) is a function definition that is not bound to an identifier.
}{function_literal}
For example, here we pass to the \texttt{map} \linkterm{function}{sml_declarations} an \linkterm{anonymous function}{function_literal} and a list:
\begin{verbatim}
- map (fn x => x * x) [2,3,4];
val it = [4,9,16] : int list
\end{verbatim}
\textbf{Note: }we can also give an \linkterm{anonymous function}{function_literal} a name (treated as \linkterm{value declaration}{sml_declarations}) and call it, e.g.:
\begin{verbatim}
- val identity = fn x => x;
val identity = fn : 'a -> 'a
- identity (5);
val it = 5 : int
\end{verbatim}
Given that \linkterm{functions}{sml_declarations} can be passed as arguments, they can be returned as well, e.g., here is a \linkterm{functions}{sml_declarations} \texttt{addconst} that takes an argument $k$ and returns another \linkterm{functions}{sml_declarations} that adds $k$ to its argument.
\begin{verbatim}
- val addconst = fn k => (fn x => x + k);
val addconst = fn : int -> int -> int
- (addconst 4) 5;
val it = 9 : int
\end{verbatim}

\definition{Cartesian vs Cascaded Function}{
We call a \linkterm{function}{sml_declarations} \textbf{cartesian}, iff the \textit{argument} type is a \textit{product} type, we call it \textbf{cascaded}, iff the \textit{result} type is a \textit{function} type. (\linkterm{uncurried}{curried_function_type} vs \linkterm{curried}{curried_function_type})
}{cartesian_cascaded}
For example, a \linkterm{cartesian}{cartesian_cascaded} plus function would look like:
\begin{verbatim}
- fun cartesian_plus (x:int, y:int) = x + y;
val cartesian_plus = fn : int * int -> int
- cartesian_plus (3,2);
val it = 5 : int
\end{verbatim}
However, a \linkterm{cascaded}{cartesian_cascaded} plus function would look like:
\begin{verbatim}
- fun cascaded_plus (x:int) = (fn y:int => x + y);
val cascaded_plus = fn : int -> int -> int
- (cascaded_plus 3) 2;
val it = 5 : int
\end{verbatim}
The following function is both \linkterm{cascaded}{cartesian_cascaded} and \linkterm{cartesian}{cartesian_cascaded}:
\begin{verbatim}
- fun both_plus (x:int, y:int) = fn (z:int) => x+y+z;
val both_plus = fn : int * int -> int -> int
- (both_plus(2,3)) 5;
val it = 10 : int
\end{verbatim}

\commandnote{
Function application always associates to the left, and function types always associate to the right. Usually we omit brackets for convenience, hence associativity must be implied.
\begin{itemize}
    \item $\operatorname{both\_plus}(2,3) 5 \leadsto (\operatorname{both\_plus}(2,3)) 5$
    \item $\operatorname{int} \to \operatorname{int} \to \operatorname{int} \leadsto \operatorname{int} \to (\operatorname{int} \to \operatorname{int})$
\end{itemize}
}



\definition{Left Folding Operator}{
\linkterm{SML}{def:sml} provides the left folding \linkterm{operator}{arithmetic} to realize a recurrent computation schema. It is a \linkterm{higher-order function}{higher_order_functions} that:
\begin{itemize}
    \item takes a \linkterm{function}{sml_declarations} \inlinecode{f : 'a * 'b -> 'b}
    \item takes a starting \linkterm{value}{sml_declarations} \inlinecode{s : 'b}
    \item takes a list \inlinecode{[x1, ... xn] : 'a list}
    \item and reduces that list into a single value by applying \inlinecode{f} from the left
\end{itemize}
}{left_folding}
The general form of the \linkterm{foldl}{left_folding} operator:
\begin{verbatim}
foldl f s [x1, x2, x3, …, xn]
= f(xn, f(x(n-1), … f(x1, s)…))
\end{verbatim}
We can think of the \linkterm{foldl}{left_folding} operator as \linkterm{recursive}{recursion} list processing. For example, remember our \texttt{sumlist} that we wrote \linkterm{recursively}{recursion}:
\begin{verbatim}
- fun sumlist nil = 0 | sumlist (h::t) = h + sumlist t;
\end{verbatim}
instead of implementing a \linkterm{recursive function}{recursion}, we can \linkterm{foldl}{left_folding} the summation operator (\texttt{op+}) and give it \texttt{s=0}, and our list to sum:
\begin{verbatim}
- foldl op+ 0 [1,2,3];
val it = 6 : int
\end{verbatim}
This is what is going on under the hood:
\begin{verbatim}
- op+ (3, op+ (2, op+ (1, 0)));
val it = 6 : int
\end{verbatim}

We also can use \linkterm{value declaration}{sml_declarations} to \linkterm{declare}{sml_declarations} a general summation function that works for any list using the form \inlinecode{val n = foldl f s} where \texttt{n} is the name that we can call. For example:
\begin{verbatim}
- val sum = foldl op+ 0;
val sum = fn : int list -> int
- sum [1,2,3,4];
val it = 10 : int
- sum [];
val it = 0 : int
\end{verbatim}

\textbf{Another Example:}
Remember the \linkterm{recursive function}{recursion} we wrote to get the length of a list:
\begin{verbatim}
- fun length nil = 0 | length (_::t) = 1 + length t;
\end{verbatim}

With \linkterm{foldl}{left_folding}, we can just write the following:
\begin{verbatim}
foldl (fn (_, acc) => acc + 1) 0 [10,20,30]
\end{verbatim}
It will evaluate to: \texttt{fn (30,fn (20,fn (10,0)))} $\leadsto$  \texttt{fn (30,fn (20,1))} $\leadsto$ \texttt{fn (30,2)} $\leadsto$ \texttt{3}.

The type of \linkterm{foldl}{left_folding} is: \texttt{foldl : ('a * 'b -> 'b) -> 'b -> 'a list -> 'b}. i.e. First argument is a \linkterm{function}{sml_declarations} \texttt{f} of type \texttt{'a * 'b -> 'b}. In the previous example, that function was \texttt{fn (\_, acc) => acc + 1}:
\begin{verbatim}
- fn (_, acc) => acc + 1;
val it = fn : 'a * int -> int
\end{verbatim} 
Second argument is a start value of type \texttt{'b}, in our case it is \texttt{0} of type \texttt{int}. The third argument is a list of type \texttt{'a list}. Finally, it returns a result of type \texttt{'b}.

\textbf{Another Example: }remember how we wrote a \linkterm{recursive function}{recursion} to reverse a list:
\begin{verbatim}
- fun reverse nil = nil | reverse (h::t) = reverse t @ [h];
val reverse = fn : 'a list -> 'a list
- reverse [1,2,3,4,5];
val it = [5,4,3,2,1] : int list
\end{verbatim}
We can instead \linkterm{foldl}{left_folding} the \texttt{op::} (cons) operator:
\begin{verbatim}
- foldl op:: nil [1,2,3,4,5];
val it = [5,4,3,2,1] : int list
\end{verbatim}
It will evaluate to: \texttt{op::(5,op::(4,op::(3,op::(2,op::(1,nil)))))} $\leadsto$
\begin{itemize}
    \item  \texttt{op::(5,op::(4,op::(3,op::(2,[1]))))} $\leadsto$
    \item \texttt{op::(5,op::(4,op::(3,[2,1])))} $\leadsto$
    \item \texttt{op::(5,op::(4,[3,2,1]))} $\leadsto$
    \item \texttt{op::(5,[4,3,2,1])} $\leadsto$
    \item \texttt{[5,4,3,2,1]}
\end{itemize}


\minititle{Check whether a number is prime}
\definition{Prime Numbers}{
An integer $p \in \mathbb{Z}$ is prime if $p>1$ and the only positive divisors of $p$ are $1$ and $p$ itself. Examples:
\begin{itemize}
    \item $2$ is prime because its divisors are $\set{1,2}$
    \item $7$ is prime because its divisors are $\set{1,7}$
    \item $9$ is not prime because its divisors are $\set{1,3,9}$
\end{itemize}
The \linkterm{infinite}{set_cardinality} \linkterm{set}{def:set} of primes is $\set{2, 3, 5, 7, 11, 13, \cdots}$
}{prime_numbers}

\definition{Composite Numbers}{
A positive number $n>1$ is called composite iff it is not prime. $1$ is neither prime nor composite by defintiion.
}{composite_numbers}

To check if a number $n$ is \linkterm{prime}{prime_numbers} we need to 
\begin{itemize}
    \item check if $n>=2$, otherwise it is $1$ which is not \linkterm{prime}{prime_numbers}.
    \item if $n=2$, then it is \linkterm{prime}{prime_numbers}.
    \item for all $n>2$ we need to see if there exists a number $d$ between $2$ and $n-1$ that divides $n$, if we find one, then $n$ is not \linkterm{prime}{prime_numbers}. If we don't find any, then $n$ is \linkterm{prime}{prime_numbers}.
\end{itemize}
For example, consider $n=5$, it does not divide any number in $[2 - 4]$ meaning the only divisors are $1$ and $5$, $\leadsto$ \linkterm{prime}{prime_numbers}. However, $6$ divides $3,2 \in [2-5] \leadsto$ not \linkterm{prime}{prime_numbers}.

To implement that \linkterm{recursively}{recursion}, we need a main \linkterm{function}{sml_declarations} (\texttt{prime n}) and a helper one (\texttt{PrimeHelper (n,d)}).  \texttt{prime n} returns \texttt{false} if $n=1$, otherwise it calls  \texttt{PrimeHelper (n,2)}. \texttt{PrimeHelper (n,d)} has two base cases:
\begin{itemize}
    \item if $n=d$ (i.e. $n=2$) $\leadsto$ \texttt{true}
    \item if $n \operatorname{mod} d = 0 \leadsto$  \texttt{false} (we found a $d$ such that, $n$ is divisible by $d$)
\end{itemize}
otherwise, we \linkterm{recursively}{recursion} call \texttt{PrimeHelper (n,d+1)}.


\begin{verbatim}
- fun PrimeHelper (n,d) =
    if d = n then true
    else if n mod d = 0 then false
    else PrimeHelper (n,d+1);

val PrimeHelper = fn : int * int -> bool

- fun prime n =
    if n<2 then false
    else PrimeHelper (n,2);

val prime = fn : int -> bool

- prime 3;
val it = true : bool
- prime 4;
val it = false : bool
- prime 9;
val it = false : bool
- prime 11;
val it = true : bool
\end{verbatim}
Let's trace:
$\operatorname{prime} 9 \leadsto \operatorname{PrimeHelper}(9, 2) \leadsto \operatorname{PrimeHelper}(9, 3) \leadsto \operatorname{false}$

\minititle{Check if a list has a prime number:}
We use our \texttt{prime n} \linkterm{function}{sml_declarations} and the \linkterm{\texttt{foldl} operator}{left_folding}. We will \linkterm{foldl}{left_folding} the list with the \linkterm{function}{sml_declarations} \texttt{fn (a,b) => (prime a) orelse b} starting with \texttt{false} as the starting value.
\begin{verbatim}
- fun foo l = foldl (fn (a,b) => (prime a) orelse b) false l;
val foo = fn : int list -> bool
- foo [4,2,6,8];
val it = true : bool
- foo [4,6,8,10];
val it = false : bool
\end{verbatim}
Let's trace:
\begin{verbatim}
foldl f false [4,2,6,8]
= f(8, f(6, f(2, f(4, false))))
= f(8, f(6, f(2, ((prime 4) orelse false))))
= f(8, f(6, f(2, false)))
= f(8, f(6, ((prime 2) orelse false)))
= f(8, f(6, true))
= f(8, ((prime 6) orelse true))
= f(8, true)
= (prime 8) orelse true
= true
\end{verbatim}



\definition{Right Folding Operator}{
\linkterm{SML}{def:sml} provides the right folding \linkterm{operator}{arithmetic} to realize a recurrent computation schema. It is a \linkterm{higher-order function}{higher_order_functions} that:
\begin{itemize}
    \item takes a \linkterm{function}{sml_declarations} \inlinecode{f : 'a * 'b -> 'b}
    \item takes a starting \linkterm{value}{sml_declarations} \inlinecode{s : 'b}
    \item takes a list \inlinecode{[x1, ... xn] : 'a list}
    \item and reduces that list into a single value by applying \inlinecode{f} from the right
\end{itemize}
}{right_folding}
The general form of the \linkterm{foldr}{right_folding} operator:
\begin{verbatim}
foldl f s [x1, x2, x3, …, xn]
= f(x1, f(x2, … f(xn, s)…))
\end{verbatim}
For example: \texttt{foldr f s [x1, x2, x3, x4] = f(x1, f(x2, f(x3, f(x4, s))))}.

A useful usecase would be to \linkterm{foldr}{right_folding} the cons operator (\texttt{op::}) which gives us an append function of two lists. We previously wrote an append function \linkterm{recursively}{recursion}:
\begin{verbatim}
- fun append (nil,y) = y | append (h::t, y) = h::append(t,y);
val append = fn : 'a list * 'a list -> 'a list
- append ([1,2],[3,4]);
val it = [1,2,3,4] : int list
\end{verbatim}
We give the \linkterm{recursive function}{recursion} a list $x$ and a list $y$ and we get back a new list wih $y$ appended to $x$. If we \linkterm{foldr}{right_folding} the cons operator, then we give first the list $y$, then $x$ to get our desired result:
\begin{verbatim}
- foldr op:: [3,4] [1,2];
val it = [1,2,3,4] : int list
\end{verbatim}
It will evaluate to: \texttt{op::(1,op::(2, [3,4]))}


\definition{Defining Equations}{
A \linkterm{set}{def:set} of equations $F(x) = \Phi_i (F,x)$ where $F$ and $x$ are variables and $\Phi_i (F,x)$ an expression in $F$ and $x$ is called defining for a \linkterm{function}{function} $f$ if $f$ solves them simultaneously for all $x$.
}{defining_equations}
For example, here is a \linkterm{defining equation system}{defining_equations} for unary addition:
\begin{verbatim}
+(n, 0)     = n
+(m, s(n))  = s(+(m, n))  
\end{verbatim}

\Theorem{
    $+ \subseteq \cartprod{\mathbb{N}_1,\mathbb{N}_1,\mathbb{N}_1}$ is a \linkterm{total function}{function}.
}{th:plus_total_function}

\Lemma{
$\forall (n,m) \in \cartprod{\mathbb{N}_1, \mathbb{N}_1}$ there is exactly one $l \in \mathbb{N}_1$ with $+(n,m)=l$.
}{lem:plus_total_function}

\Proof{
By \linkterm{induction}{induction} on $m$. We have two cases:

1. Base case $(m=0)$
\begin{itemize}
    \item Choose $l:=n$, so we have $+(n,0) = n = l$
    \item For any $l' = +(n,0)$ we have $l' = n = l$ 
\end{itemize}
2. \linkterm{Induction step}{induction_step} ($m = s(k)$)
\begin{itemize}
    \item Assume that there is a unique $r = +(n,k)$, choose $l:=s(r)$, so we have $+(n,s(k)) = s(+(n,k)) = s(r)$
    \item Again, for any $l' = +(n,s(k))$ we have $l' = l$
\end{itemize}
}
\Corollary{
$+: \cartprod{\mathbb{N}_1, \mathbb{N}_1} \to \mathbb{N}_1$ is a \linkterm{total function}{function}
}{cor:plus_total_function}

\textbf{More Examples:}
\begin{itemize}
    \item the equations: $\delta(0) = 0$ and $\delta(s(n)) = s(s(\delta(n)))$ are \linkterm{defining equations}{defining_equations} for the doubling \linkterm{function}{function}.
    \item the equations: $\mu(l,0)=0$ and $\mu(l,s(r)) = +(\mu(l,r), l)$ are \linkterm{defining equations}{defining_equations} for the \linkterm{multiplication}{multiplication} \linkterm{function}{function}.
    \item the equations; $i(a,A,B)=A$ and $i(b,A,B)=B$ where $A,B$ are \linkterm{unary natural number}{unarynatnums} are \linkterm{defining equations}{defining_equations} for the \texttt{if-then-else} \linkterm{function}{function} on \linkterm{unary natural number}{unarynatnums} where $a$ represents \textit{true} and $b$ is \textit{false}.
    \item the equations: $\mu(0) = 0$, $\mu(s(0)) = 0$, and $\mu(s(s(n))) = s(n)$ are \linkterm{defining equations}{defining_equations} for the \linkterm{function}{function} $\mu$ on \linkterm{unary natural number}{unarynatnums} where $\mu$ is the \texttt{lower halving} \linkterm{function}{function} $\mu:n \mapsto \lfloor \frac{n}{2} \rfloor $.
\end{itemize}

\definition{Inductively Defined Sets}{
An inductively defined set $\tuple{S,C}$ is a \linkterm{set}{def:set} $S$ together with a \linkterm{finite}{set_cardinality} \linkterm{set}{def:set} $C:= \setbuilder{c_i}{1 \leq i \leq n}$ of $k_i$-ary constructors $c_i: S^{k_i} \to S$ with $k_i \geq 0$, such that:
\begin{itemize}
    \item if $s_i \in S \quad \forall 1 \leq i \leq k_i$, then $c_i(s_1, \cdots, s_{k_i}) \in S$ (\textit{generated by constructors})
    \item all constructors are \linkterm{injective}{injective} (\textit{no internal confusion})
    \item $\intersection{\operatorname{Im}(c_i), \operatorname{Im}(c_j)} = \phi$ for $i \neq j$ (\textit{no inter-constructor confusion})
    \item for every $s \in S$ there is a constructor $c \in C$ with $s \in \operatorname{Im}(c)$ (\textit{no junk})
    \item If $k_i = 0$, it is a nullary constructor (a constant element).
\end{itemize}
}{inductively_defined_set}
\textbf{Example: }$\tuple{\mathbb{N}_1, \set{s,0}}$ is an \linkterm{inductively defined set}{inductively_defined_set}.
\Proof{

1. Generated by constructors:

This is guranteed by $P1$ and $P2$. $0 \in \mathbb{N}_1$ and if $n \in \mathbb{N}_1$ then $s(n) \in \mathbb{N}_1$.

2. No internal confusion: $c_i (s_1, \cdots, s_{k_i}) = c_i(t_1, \cdots, t_{k_i}) \Rightarrow s_j = t_j \forall j$.

$P4$ gurantees this. If $s(n) = s(m)$ then $n=m$

3. No inter-constructor confusion:

$P3$ asserts that. Different constructors never produce the same element. The element $0$ is not the image $s$, nothing is both a zero and a successor.

4. No junk is assured by $P5$: $\forall s \in S, \exists c \in C \text{ such that }s \in \operatorname{Im}(c)$. i.e. every element of $S$ comes from some constructor. There are no extra elements that didn't come from constructors.
}

We can define Lists of natural numbers $\mathcal{L}[\mathbb{N}]$ \linkterm{inductively}{inductively_defined_set} as well. Let the \texttt{nil}-rule be starting with the empty list $[]$, and the \texttt{cons}-rule extends the list by adding a number $n \in \mathbb{N}$ at the front. For example $\operatorname{cons}(3,\operatorname{cons}(2, \operatorname{cons}(1, \operatorname{nil}))) = [3,2,1]$ and $[] = \operatorname{nil}$

\definition{List Piano Axioms}{
We call the following \linkterm{set}{def:set} of \linkterm{axioms}{axiom} \textit{list Peano axioms} for $\mathcal{L}[\mathbb{N}]$ in analogy to the \linkterm{Peano Axioms}{piano} in \Defref{piano}.
\begin{itemize}
\item \textbf{LP1}: $\operatorname{nil} \in \mathcal{L}[\mathbb{N}]$.
\item \textbf{LP2}: $\operatorname{cons}: \mathbb{N} \times \mathcal{L}[\mathbb{N}] \to \mathcal{L}[\mathbb{N}]$.
\item \textbf{LP3}: \texttt{nil} is not a \texttt{cons}-value.
\item \textbf{LP4}: \texttt{cons} is \linkterm{injective}{injective}.
\item \textbf{LP5}:
\begin{enumerate}
\item \textbf{Base Case}: \texttt{nil} has property $P$. ($P(\operatorname{nil})$)
\item \textbf{Inductive Step}: If a list $l$ has property $P$, then for every natural number $n \in \mathbb{N}$, the list $\operatorname{cons}(n,l)$ also has property $P$. i.e.
\[
P(l) \Rightarrow P(\operatorname{cons}(n,l)) \quad \forall n \in \mathbb{N}, \forall l \in \mathcal{L}[\mathbb{N}]
\]
\item \textbf{Conclusion}: $P(l) \quad \forall l \in \mathcal{L}[\mathbb{N}]$
\end{enumerate}
\[
\left( P(\operatorname{nil}) \land \forall l \in \mathcal{L}[\mathbb{N}], \forall n \in \mathbb{N}, (P(l) \Rightarrow P(\operatorname{cons}(n,l))) \right) \Rightarrow \forall l \in \mathcal{L}[\mathbb{N}], P(l)
\]
\end{itemize}
}{piano_list}

\textbf{Example (Length of a List).}  
We define the length function on lists as follows:
\[
\text{len}(nil) = 0, \qquad \text{len}(\text{cons}(n,l)) = 1 + \text{len}(l).
\]

We want to prove
\[
\forall l \in \mathcal{L}[\mathbb{N}], \; \text{len}(l) \in \mathbb{N}.
\]

\Proof{
We proceed by list induction (Axiom LP5).

\emph{Base Case:} Let $l = nil$. Then
\[
\text{len}(nil) = 0 \in \mathbb{N}.
\]
Thus the base case holds.

\emph{Inductive Step:} Assume $\text{len}(l) \in \mathbb{N}$ for some $l \in \mathcal{L}[\mathbb{N}]$ (induction hypothesis).  
We need to show that $\text{len}(\text{cons}(n,l)) \in \mathbb{N}$ for any $n \in \mathbb{N}$.  
Indeed,
\[
\text{len}(\text{cons}(n,l)) = 1 + \text{len}(l).
\]
By the induction hypothesis, $\text{len}(l) \in \mathbb{N}$. Since $\mathbb{N}$ is closed under addition, we have $1 + \text{len}(l) \in \mathbb{N}$.

\emph{Conclusion:} By the principle of list induction, it follows that
\[
\forall l \in \mathcal{L}[\mathbb{N}], \quad \text{len}(l) \in \mathbb{N}.
\]
}


\textbf{Example (Double Reversal).}  
We define the reverse function on lists as follows:
\[
\begin{aligned}
\text{rev}(nil) &= nil, \\
\text{rev}(\text{cons}(n,l)) &= \text{append}(\text{rev}(l), \text{cons}(n,nil)),
\end{aligned}
\]
where $\text{append}$ is the usual concatenation of two lists.

We want to prove
\[
\forall l \in \mathcal{L}[\mathbb{N}], \quad \text{rev}(\text{rev}(l)) = l.
\]

\Proof{
We proceed by list induction (Axiom LP5).

\emph{Base Case:} Let $l = nil$. Then
\[
\text{rev}(\text{rev}(nil)) = \text{rev}(nil) = nil,
\]
so the base case holds.

\emph{Inductive Step:} Assume $\text{rev}(\text{rev}(l)) = l$ for some $l \in \mathcal{L}[\mathbb{N}]$ (induction hypothesis).  
We need to show that
\[
\text{rev}(\text{rev}(\text{cons}(n,l))) = \text{cons}(n,l), \qquad \forall n \in \mathbb{N}.
\]

We compute:
\[
\begin{aligned}
\text{rev}(\text{rev}(\text{cons}(n,l))) 
&= \text{rev}(\text{rev}(\text{append}(\text{rev}(l), \text{cons}(n,nil)))) \\
&= \text{rev}\Big(\text{append}(\text{rev}(\text{cons}(n,nil)), \text{rev}(\text{rev}(l)))\Big).
\end{aligned}
\]

By the induction hypothesis, $\text{rev}(\text{rev}(l)) = l$.  
Also, $\text{rev}(\text{cons}(n,nil)) = \text{append}(\text{rev}(nil), \text{cons}(n,nil)) = \text{cons}(n,nil)$.  
So:
\[
\text{rev}(\text{rev}(\text{cons}(n,l))) = \text{rev}(\text{append}(\text{cons}(n,nil), l)).
\]

Now apply the auxiliary lemma that
\[
\text{rev}(\text{append}(x,y)) = \text{append}(\text{rev}(y), \text{rev}(x)).
\]
Using this with $x = \text{cons}(n,nil)$ and $y = l$, we get:
\[
\text{rev}(\text{append}(\text{cons}(n,nil), l)) = \text{append}(\text{rev}(l), \text{rev}(\text{cons}(n,nil))).
\]

But $\text{rev}(\text{cons}(n,nil)) = \text{cons}(n,nil)$, hence
\[
\text{rev}(\text{rev}(\text{cons}(n,l))) = \text{append}(l, \text{cons}(n,nil)) = \text{cons}(n,l).
\]

Thus the inductive step is proved.

\emph{Conclusion:} By the principle of list induction, it follows that
\[
\forall l \in \mathcal{L}[\mathbb{N}], \quad \text{rev}(\text{rev}(l)) = l.
\]
}

\definition{Datatype Declaration}{
\linkterm{SML}{def:sml} datatype declarations provide a concerete syntax for \linkterm{inductively defined sets}{inductively_defined_set} (called (abstract) data types in programming) via the keyword \texttt{datatype} followed by a sequence of \textbf{constructor declarations}
}{datatype_declaration}

\textbf{Example: }We can declare a datatype for \linkterm{unary natural numbers}{unarynatnums} by:
\begin{verbatim}
- datatype mynat = zero | suc of mynat;
datatype mynat = suc of mynat | zero

- val n0 = zero;
val n0 = zero : mynat

- val n1 = suc zero;
val n1 = suc zero : mynat

- val n2 = suc (suc zero);
val n2 = suc (suc zero) : mynat

- val n3 = suc (suc (suc zero));
val n3 = suc (suc (suc zero)) : mynat
\end{verbatim}

Therefore, we can implement \linkterm{functions}{sml_declarations} on our datatype \texttt{mynat}, e.g. here is unary addition:
\begin{verbatim}
- fun add (zero, m) = m | add (suc n,m) = suc (add (n,m));
val add = fn : mynat * mynat -> mynat

- add (suc zero, suc (suc zero));
val it = suc (suc (suc zero)) : mynat
\end{verbatim}

\minititle{Converting unary to int:}
\begin{verbatim}
- fun num (zero) = 0 | num (suc n) = num (n) + 1;
val num = fn : mynat -> int

- num (suc (suc (suc zero)));
val it = 3 : int
\end{verbatim}

\minititle{Check even-odd unary using Mutual Recursion:}
\begin{verbatim}
- fun even (zero) = true | even (suc n) = odd (n)
      and
      odd (zero) = false | odd (suc n) = even (n);

val even = fn : mynat -> bool
val odd = fn : mynat -> bool

- odd (suc (suc (suc zero)));
val it = true : bool

- even (suc (suc (suc (suc zero))));
val it = true : bool
\end{verbatim}
Trace:
\begin{verbatim}
even (suc (suc (suc (suc zero)))) = odd (suc (suc (suc zero))) = even(suc (suc zero))
= odd (suc zero) = even (zero) = true
\end{verbatim}

We also can declare \textbf{Enumeration Types}, e.g. here is a type for weekdays:
\begin{verbatim}
- datatype day = mon | tue | wed | thu | fri | sat | sun;
datatype day = fri | mon | sat | sun | thu | tue | wed
\end{verbatim}
And here is a function that check if a given weekday is weekend or not:
\begin{verbatim}
- fun weekend sat = true | weekend sun = true | weekend _ = false;
val weekend = fn : day -> bool

- weekend fri;
val it = false : bool

- weekend sat;
val it = true : bool
\end{verbatim}

\minititle{Binary Tree datatype (only leafs are valued):}
\begin{verbatim}
datatype tree = leaf of int | node of tree * tree;

- val t = node (leaf 1, leaf 2);
val t = node (leaf 1,leaf 2) : tree
\end{verbatim}

\minititle{Binary Tree datatype (every node is valued):}
\begin{verbatim}
- datatype tree = leaf of int | node of int * tree * tree;

- val t1 = node (1, leaf 2, leaf 3);
val t1 = node (1,leaf 2,leaf 3) : tree
\end{verbatim}

\minititle{A function to sum node values in a binary tree:}
\begin{verbatim}
- fun sum (leaf n) = n
    | sum (node (v, l, r)) = v + sum l + sum r;

val sum = fn : tree -> int

- sum t1;
val it = 6 : int
\end{verbatim}

\minititle{Number of leaves:}
\begin{verbatim}
- datatype tree = leaf of int | node of int * tree * tree;

- fun leavescount (leaf x) = 1 | leavescount (node (l, r)) = leavescount l + leavescount r;
val leavescount = fn : tree -> int

- val t1 = node (leaf 1, node (leaf 2, leaf 3));

- leavescount t1;
val it = 3 : int
\end{verbatim}

\minititle{Multiway tree:}
\begin{verbatim}
datatype t = a of int | c of t list

- val v = a 1;
val v = a 1 : t

- val v = c [a 1, c [a 2]];
val v = c [a 1,c [a 2]] : t
\end{verbatim}


\minititle{Sum all nodes in the tree:}
\begin{verbatim}
- fun sum (a n) = n
    | sum (c children) = foldl (fn (sub, acc) => sum sub + acc) 0 children;

val sum = fn : t -> int

- sum v;
val it = 3 : int
\end{verbatim}


\textbf{Another Example:} We can declare types for geometric shapes:
\begin{itemize}
    \item A circle is defined by only one value $r \in \mathbb{R}^{+}$ (its radius)
    \item A square is defined by one value $a \in \mathbb{R}^{+}$ (side length)
    \item A triangle is defined by the three values $a,b,c \in \cartprod{\mathbb{R}^{+}, \mathbb{R}^{+}, \mathbb{R}^{+}}$ (two sides and base)
\end{itemize}

\begin{verbatim}
- datatype shape =
      Circle of real
    | Square of real
    | Triangle of real * real * real;

- Circle 4.0;
val it = Circle 4 : shape

- Square 3.0;
val it = Square 3 : shape

- Triangle (4.0,3.0,5.0);
val it = Triangle (4,3,5) : shape
\end{verbatim}
And here is a function that compute the area of a shape:
\begin{verbatim}
- fun area (Circle r) = Math.pi * r * r
    | area (Square a) = a * a
    | area (Triangle (a,b,c)) =
        let
            val s = (a + b + c)/2.0
        in
            Math.sqrt(s * (s-a) * (s-b) * (s-c))
        end;

val area = fn : shape -> real

- area (Square 3.0);
val it = 9 : real

- area (Circle 4.0);
val it = 50.2654824574 : real

- area (Triangle (4.0, 3.0, 5.0));
val it = 6 : real
\end{verbatim}

Until now, our \linkterm{functions}{sml_declarations} have been characterized entirely by their \linkterm{values}{sml_declarations} on their arguments (as \linkterm{mathematical}{mathematics} \linkterm{functions}{function} behave). However, \linkterm{SML}{def:sml} also considers effects, e.g. for
\begin{itemize}
    \item \textbf{input/output}: the interesting bit about a \texttt{print} statement is the effect,
    \item \refterm{mutation}{mutation_sml}: allocation and modification of storage during evaluation,
    \item \refterm{communication}{communication_sml}: data may be send and received over channels, and
    \item \textbf{exceptions}: abort evaluation by signaling an exceptional condition.
\end{itemize} 
\commandnote{
    An effect is any action resulting from an evaluation that is not returning a value.
}

Input and Output is handled via \textit{streams}. There are two predefined streams \texttt{TextIO.stdIn} and \texttt{TextIO.stdOut}.

\textbf{Example: } Input via \texttt{TextIO.inputLine}: \texttt{TextIO.instream} $\to$ \texttt{string}
\begin{verbatim}
- TextIO.inputLine(TextIO.stdIn);
ibrahim
val it = SOME "ibrahim\n" : string option
\end{verbatim}

\textbf{Example: } Printing to Standard Output. \texttt{TextIO.print} prints its arguments to \texttt{stdIn}. \texttt{print} also returns a value \texttt{()} of type \texttt{unit}:
\begin{verbatim}
- print "Hello World! \n";
Hello World!
val it = () : unit
\end{verbatim}
\commandnote{
Users can create streams as files using \texttt{TextIO.openIn} and \texttt{TextIO.openOut}. However, streams should be closed when no longer needed via \texttt{TextIO.closeIn} and \texttt{TextIO.closeOut}. 
}

\definition{Exception}{
An \textbf{exception} is a special \linkterm{SML}{def:sml} object. \textbf{Raising} an exception $e$ in a \linkterm{function}{sml_declarations} aborts functional computation and returns $e$ to the next level.
}{exception_sml}
\textbf{Example: }Predefined \linkterm{exceptions}{exception_sml}:
\begin{verbatim}
- 3 div 0;
uncaught exception Div [divide by zero]
  raised at: stdIn:88.3-88.6

- fib 100;
uncaught exception Overflow [overflow]
  raised at: <file stdIn>
\end{verbatim}

\textbf{Example: }User-defined \linkterm{exceptions}{exception_sml}:
\begin{verbatim}
- exception Empty;
exception Empty

- Empty;
val it = Empty(-) : exn
\end{verbatim}

\textbf{Example: } \linkterm{Exception}{exception_sml} constructors:
\begin{verbatim}
- exception SysError of int;
exception SysError of int

- SysError;
val it = fn : int -> exn
\end{verbatim}

\textbf{Example: }Programming with \linkterm{exceptions}{exception_sml} (factorial):
\begin{verbatim}
- exception Factorial;
exception Factorial

- local
    fun fact 0 = 1 | fact n = n * fact (n-1)
  in  
    fun safe_factorial n =
        if n >= 0 then fact n else raise Factorial
  end;

val safe_factorial = fn : int -> int

- safe_factorial(~1);
uncaught exception Factorial
  raised at: stdIn:112.34-112.43
\end{verbatim}

\commandnote{
\texttt{local} act just like \texttt{let}, only that is also that the body (between \texttt{in} and \texttt{end}) contains sequence of \linkterm{declarations}{sml_declarations} and not an expression whose value is returned as the value of the overall expression. Here the identifier \texttt{fact} is bound to the \linkterm{recursive function}{recursion} in the definition of \texttt{safe\_factorial} but unbound outside. This way we avoid polluting the name space with \linkterm{functions}{sml_declarations} that are only used \linkterm{locally}{local_declaration_sml}.
}

\minititle{Exception Handling:}
\begin{verbatim}
 exception NaN;
 fun read_integer () =
    let
        val intstring = case TextIO.inputLine(TextIO.stdIn) of
                           NONE => raise NaN
                         | SOME(s) => s
    in
        case Int.fromString intstring of
             NONE => raise NaN
           | SOME i => i
    end;   
\end{verbatim}
\begin{verbatim}
fun factorial_driver () =
    let val input = read_integer ()
        val result = Int.toString (safe_factorial input)
    in
        print result
    end
    handle Factorial => print "Out of range.\n"
         | NaN => print "Not a Number!\n";
\end{verbatim}

\begin{verbatim}
- factorial_driver ();
3
6
val it = () : unit

- factorial_driver ();
-3
Out of range.
val it = () : unit

- factorial_driver ();
hi
Not a Number!
val it = () : unit
\end{verbatim}